\documentclass[a4paper,10pt]{report}
\usepackage[utf8]{inputenc}
\usepackage[pdftex]{graphicx}
\usepackage{fancyhdr}
\usepackage{listings}
\usepackage[francais]{babel}
\usepackage{hyperref}
\usepackage{amsmath}

% Title Page
\title{IA01 : TP2}
\author{Par Camille Gerin-Roze et Thomas Perrin}



\renewcommand{\footrulewidth}{1pt}
\fancyfoot[C]{\textbf{\thepage}}
\fancyfoot[L]{}

\hypersetup{
	colorlinks=true,
	linkcolor=black
}

\begin{document}

\begin{titlepage}

\begin{flushright}
  \includegraphics[scale = 0.2]{logo_utc.jpg}
\end{flushright}
\vspace*{5cm}

\newcommand{\HRule}{\rule{\linewidth}{0.5mm}} % Defines a new command for the horizontal lines, change thickness here
\center % Center everything on the page



%----------------------------------------------------------------------------------------
%	TITLE SECTION
%----------------------------------------------------------------------------------------

\HRule \\[0.4cm]
{ \LARGE \bfseries  IA01 : Rapport du TP3}\\[0.4cm] % Title of your document
\HRule \\[1.5cm]

%----------------------------------------------------------------------------------------
%	AUTHOR SECTION
%----------------------------------------------------------------------------------------
\begin{minipage}{0.4\textwidth}
\begin{flushleft} \large
\emph{Auteurs:}\\
Camille \textsc{Gerin-Roze} \newline
Thomas \textsc{Perrin}
\end{flushleft}
\end{minipage}
~
\begin{minipage}{0.4\textwidth}

\end{minipage}\\[1.3cm]


{\large Le 3 Janvier 2016} % Date, change the \today to a set date if you want to be precise

\end{titlepage}
\tableofcontents
\chapter*{Introduction}
\chapter{Connaissances nécessaires au système expert}
 \section{Paramètres à fournir au système}
 Afin de déterminer le nombre de nains nécessaire à un chantier, plusieurs paramètres seront nécessaire :
\begin{itemize}
\item La longueur du tunnel
\item Sa hauteur
\item Sa largeur
\item Le nombre de jours maximum
\item Le type de roche
\item Le type de pioche employé
\end{itemize}

Avec ces paramètres, on va pouvoir déterminer si oui ou non le chantier est réalisable. Si il l’est, on obtiendra un nombre total de nains avec les différents rôles de chacun. Un indicateur sera présent pour dire si une équipe de nuit est nécessaire ou non pour compléter le chantier en temps et en heures.

\section{Traduction des connaissances en règles}

La base de règles est décomposé en plusieurs parties :
\begin{itemize}
\item Les règles du type de roche (RTR)
Ces règles vont permettre, en fonction du type de roche, de fixer la vitesse de travail d’un nain par jour.

\item Les règles du type de pioche (RTP)
Ces règles vont ajouter un coefficient multiplicateur à la vitesse de base, pour aller plus ou moins vite en fonction des moyens de l’employeur.

\item Les règles concernant la réalisation du chantier (RCR)
Ces règles vont déterminer si le chantier est réalisable avec les paramètres énoncés. En effet, en fonction de la vitesse des nains, de la hauteur, largeur et longueur du tunnel, on peut d’avance vérifier si cela tiendra avec le nombre de jours énoncé ou non.

\item Les règles calculant le nombre de nains (RN)
Ces dernières règles vont conclure en calculant le nombre de nains totals.
\end{itemize}

Toutes les règles ont été déterminés par la vidéo présentés ci-dessus, et par des échanges avec John Lang, le créateur de l’univers du Donjon de Naheulbeuk. Nous avons seulement modifié une seule règle pour la rendre cohérente, celle des nains de ravitaillements en bière. 



\chapter{Implémentation du système expert}
  \section{Choix de la représentation Lisp}
    \subsection{Base de règles}
    Pour représenter les règles de notre SE, nous avons choisi d'utiliser une liste qui contiendra 2 sous-listes ( les nouveaux faits et les prémisses )
    ainsi que le nom de la règle. Par exemple, voici la règle RCR1 qui permet de savoir si un chantier est réalisable:\newline
    \begin{lstlisting}[language=Lisp]
(
	( ;;;; DEBUT DES NOUVEAUX FAITS
		(ChantierRealisable . T)
		(EquipeDeNuit . NIL)
	);;;; FIN DES NOUVEAUX FAITS
		(;;;;DEBUT DES PREMISSES
		(COMPARAISON (>=
			(* LargeurTunnel HauteurTunnel
			    NombreDeJours VitesseNain)
			LongueurTunnel
		))
	);;;;FIN DES PREMISSES
RCR1);;;;NOM DE LA REGLE

    \end{lstlisting}

    On peut noter que les premisses ont un format bien spécial. En effet, chaque condition commence par le type de vérification à faire. Cela peut être une égalité,
    une comparaison, ou encore une définition (savoir si un fait a déjà été défini). Cette étiquette permet de simplifier le traitement lors de la vérification
    des prémisses. Pour les prémisses, une fonction va donc vérifier si LargeurTunnel * HauteurTunnel * NombreDeJours * VitesseNain est bien supérieur ou égale à la
    longueur du tunnel. Si c'est le cas, notre moteur d'inférence va apeller une nouvelle fonction qui va évaluer les valeurs des nouveaux faits et les ajouter à
    notre base de fait. Nous avons également une A-List qui contient la description de chaque règles, mais celle-ci a été mise à part dans le but de garder un code
    clair. Cette représentation permet non seulement d'avoir un format homogène pour toutes les règles, et donc de traiter chaque règle de la même façon, mais elle
    permet également de traiter des conditions complexes qui demandent l'évaluation de plusieurs opérations ou encore des opérations qui demandent de rechercher
    dans la base de faits.

    \subsection{Base de faits}
      La base de faits a été faites de la manière la plus simple possible, celle-ci est juste une liste donc chaque membre est une sous-liste qui représente un fait.
      Un fait est tout simplement constitué d'une étiquette, donc le nom du fait, et d'une valeur qui peut être un nombre, un symbole, ou même une valeur booléenne.
      \begin{lstlisting}[language=Lisp]
(
  (VitesseNain 3)
  (TypeRoche GRANITE)
  (LargeurTunnel 3)
  (LongueurTunnel 28)
  (HauteurTunnel 1)
)
      \end{lstlisting}
      Ici on peut voir qu'il y a 4 faits dans notre base de faits : la vitesse des nains, le type de roche, la largeur et la longueur du tunnel, et la hauteur du tunnel.
  \section{Le moteur d'inférences}
  Le moteur d'inférence est un moteur à chaînage avant avec un fonctionnement simple, mais pour bien comprendre comment il fonctionne, il faut expliciter certaines
  fonctions qui permettent de savoir si les prémisses sont satisfaits ou bien d'ajouter les faits à notre base de faits. Nous allons donc commencer par présenter
  toutes ces fonctions.\newpage
    \subsection{Validation des prémisses}
    Pour que les prémisses soient validés, il faut que toutes les conditions soient validées. Donc nous allons itérer sur toutes les conditions de la liste
    des prémisses pour vérifier qu'elles sont toutes validée. Mais certaines conditions sont complexes et demandent d'aller vérifier l'existence ou la valeur
    de faits dans la base de faits. Pour cela nous avons défini une fonction evalOperande qui va nous permettre d'évaluer la valeur de chaque opérande :
    \begin{lstlisting}[language=Lisp]
(defun evalOperande(operande)
  (let ((newList (list)))
      (dolist (x operande newList)
        (cond
          ((LISTP x)
            (setq newList (append newList (list (evalOperande x)))))
          ((OR (EQUAL x '*) (EQUAL x '/) (EQUAL x '-) (EQUAL x '+))
            (setq newList (APPEND newList (list x))))
          ((SYMBOLP x)
            (if (NULL (ASSOC x *BaseFaits*))
	      ;;;; SI LE SYMBOLE N'EST PAS DEFINI DANS LA BDF
                (return-from evalOperande NIL) 
              (setq newList (APPEND newList (cdr (ASSOC x *BaseFaits*))))
            ))
          (T (setq newList (APPEND newList (list x))))
        )
    )
  )
)    
    \end{lstlisting}

    Pour évaluer un opérande, on va vérifier si celui-ci est une liste, si c'est une liste on va alors appeller notre fonction récursivement sur chaque éléments
    de la liste. Dans le Cas où c'est un symbole, on va alors aller chercher sa valeur dans la base de faits. Si jamais la valeur n'est pas dans la base de faits,
    il est important que la fonction renvoie NIL pour qu'on sache que cela ne peut pas être évalué et que la condition n'est pas respectée. À la fin de la fonction,
    on aura une liste qui pourra être évaluée grâce à un simple eval.\newline
    Ensuite il y a deux autres fonctions, une qui va vérifier si une condition est respectée, et une qui vérifiera si toutes les conditions sont vérifiées. La fonction
    qui vérifie si une condition est respectée va d'abbord vérifier l'étiquette placée au début de celle-ci (COMPARAISON, DEFINI, EGALITE) pour savoir comment elle 
    doit la traiter.\newpage
		\begin{lstlisting}[basicstyle = \footnotesize, language=Lisp]
(defun conditionRespecte?(condition)
(cond
  ((EQUAL (car condition) 'EGALITE)
      (return-from conditionRespecte? 
	  (EQUAL (cadr (ASSOC (car (cadr condition)) *BaseFaits*)) 
	   (cadr (cadr condition)))))
  ((EQUAL (car condition) 'DEFINI)
      (return-from conditionRespecte? 
	  (NOT (NULL (ASSOC (car (cadr condition)) *BaseFaits*)))))
  ((EQUAL (car condition) 'COMPARAISON)
    (return-from conditionRespecte?
      (eval
        (let ((expression))
          (dolist (x (cadr condition) expression)
            (cond
              ((estUnComparateur? x)
                (setq expression (append expression (list x))))
                ((listp x)
                  (if (evalOperande x)
                    (setq expression 
		      (append expression (list (eval (evalOperande x))))
		    )
                      NIL
                    ))
                (T
                  (if (ASSOC x *BaseFaits*)
                    (setq expression 
		      (append expression (list (cadr (ASSOC x *BaseFaits*))))
		     )
                    NIL
                    ))
            )
          )
        ))))
  )
)		


(defun premisseRespecte?(regle)
  (let ((test T))
    (dolist ( x (getPremisse regle) test)
      (setq test (AND test (conditionRespecte? x)))
    )
  )
)
		\end{lstlisting}

Un fois que la condition est vérifiée, on peut tout simplement faire une boucle sur toutes les conditions pour savoir si celles-ci sont toutes vrais, si elles sont 
toutes vrais on pourra ajouter les nouveaux faits à notre base de fait.\newpage
    \subsection{Calcul des nouvelles valeurs}
    \subsection{Moteur à chaînage avant}
    \subsection{Test du moteur d'inférences}
\chapter{Scénarios d'utilisation}
\chapter*{Conclusion}

  \end{document}
